\documentclass{article}
\usepackage[utf8]{inputenc}
\usepackage[T1]{fontenc}
\usepackage{geometry}
\geometry{
    a4paper,
    left=2.5cm,
    right=2.5cm,
    top=2.5cm,
    bottom=2.5cm
}
\usepackage{graphicx}
\usepackage{float}
\usepackage{hyperref}
\usepackage[english]{babel}
\usepackage{amsmath}
\usepackage{amssymb}
\usepackage{array}
\usepackage{booktabs}
\usepackage{caption}
\usepackage{subcaption}
\usepackage{enumitem}
\usepackage{fancyhdr}
\usepackage{titlesec}
\usepackage{url}
\usepackage{textcomp}
\usepackage{xcolor}
\usepackage{listings}
\usepackage{multirow}
\usepackage{longtable}
\usepackage{tabularx}
\usepackage{times}

% Set up page style
\pagestyle{fancy}
\fancyhf{}
\rfoot{\thepage}
\renewcommand{\headrulewidth}{0pt}

% Title formatting
\titleformat{\chapter}[display]{\normalfont\bfseries}{\chaptertitle\thechapter}{0pt}{\Huge}
\titlespacing*{\chapter}{0pt}{50pt}{40pt}

% Document info
\title{Template Document}
\author{Academic Support Agent}
\date{\today}

\begin{document}
\maketitle

\maketitle

\abstract{This is an example of abstract for a submission. It should be replaced by an abstract concisely describing the submission's content.}

\section{Introduction}
This template has been created by the Microservices Community \cite{MicroservicesCommunity}, with the aim of templating contributions for the \textit{International Conference on Microservices}. 
It can be loaded in a \LaTeX{} document by setting
\begin{center}
    \verb|\documentclass{mc-abstract}|
\end{center}
in the document preamble.
We hereby illustrate an example of filled template, to also provide guidelines on how to fill it.

\section{Title Page}
The title page is automatically generated by the initial \verb|\maketitle| command. For the command to work, please provide the following information in the preamble of the main \LaTeX{} file: 
\begin{itemize}
    \item \verb|\year| to set the conference year (displayed in the box at the top), e.g., to \verb|2023| as shown in the current example,
    \item \verb|\title| to set the title of the contribution,
    \item \verb|\author| to list the names of the authors of the contribution, separated by \verb|\and|, 
    \item \verb|\affiliation| to provide the affiliation of listed authors, and
    \item \verb|\correspondence| to provide an email address for correspondence on the contribution.
\end{itemize}
Please use Math apexes (like, e.g., \verb|$^1$| and \verb|$^2$|) to associate authors to affiliations, in the cases when different authors have different affiliations.

The \verb|\maketitle| command should then be followed by the \verb|\abstract| command, to which the submission abstract should be provided as parameter. For instance, the command 
\begin{center}
\verb|\abstract{This is an example of abstract}| 
\end{center}
will set the abstract to \enquote{This is an example of abstract}.
    

\section{Using \LaTeX{} Environments}
\label{sec:examples}
This section illustrates how to use \LaTeX{} environments, from subsections and paragraphs, to figures, tables, and formulas.

\subsection{Examples of Figures and Tables}
This is an example of subsection, in which we provide examples of usage of \LaTeX{} environments.
Figures and tables can actually be included with the classical \texttt{figure} and \texttt{table} environments.
An example of figure is given in \Cref{fig:example}, whereas an example of table is given in \Cref{tab:example}.

\begin{figure}
    \centering
    \includegraphics[width=2cm]{fig.png}
    \caption{Example of figure.}
    \label{fig:example}
\end{figure}
\begin{table}
    \centering
    \begin{tabular}{|rccccc|}
         \hline 
         & \textbf{1st column} & \textbf{2nd column} & \textbf{3rd column} & \textbf{4th column} & \textbf{5th column} \\
         \hline
         \textbf{1st Row} & Content & Content & Content & Content & Content \\
         \textbf{2nd Row} & Content & Content & Content & Content & Content \\ 
         \textbf{3rd Row} & Content & Content & Content & Content & Content \\ 
         \hline
    \end{tabular}
    \caption{Example of table.}
    \label{tab:example}
\end{table}

\subsection{Examples of Formulas}
This is another example of subsection, in which we provide examples of paragraphs with formulas.

\paragraph{Adding Formulas}
A formula can be added with classical \LaTeX{} \enquote{math} environments, like in the following quite popular example: 
$$ e = mc^2 $$
Formulas can be even more complex, like using \texttt{array}s as shown hereafter:
\[
\left\langle
\begin{array}{ccc}
a & b & c \\ d & e & f \\ g & h & i 
\end{array}
\right\rangle
\qquad
\left\uparrow
\begin{array}{ccc}
a & b & c \\
d & e & f \\
g & h & i 
\end{array}
\right\}
\]

\section{Conclusions}
That's it. Please suitably fill the template and send your contribution to the upcoming edition of the \textit{International Conference on Microservices}! \Smiley
\end{document}